%% Chapter/chapter4.tex
%% บทที่ 4 - ผลการทดลองและอภิปรายผล

\chapter{ผลการทดลองและอภิปรายผล}
% (English: Experimental Results and Discussion)

\noindent\hspace{\parindent}%
[กล่าวนำถึงผลการทดลองที่จะนำเสนอในบทนี้ อธิบายภาพรวมของผลการทดลอง]

%%------------------------------------------------------------------
\section{ผลการทดลองส่วนที่ 1}

\subsection{ชื่อหัวข้อย่อย}

\subsubsection{กรณีศึกษาที่ 1}

\noindent\hspace{\parindent}%
[อธิบายผลการทดลองในกรณีศึกษาที่ 1 โดยอ้างอิงถึงภาพและตารางที่แสดงผล]

% \begin{figure}[htbp]
%   \centering
%   % ใช้ subcaption สำหรับภาพย่อยหลาย ๆ ภาพ
%   \begin{subcaptiongroup}
%     \includegraphics[width=0.45\textwidth]{Figure/result_f1_case1.png}
%     \phantomcaption\label{fig:f1_case1}
%     \includegraphics[width=0.45\textwidth]{Figure/result_f2_case1.png}
%     \phantomcaption\label{fig:f2_case1}
%   \end{subcaptiongroup}
%   \caption{[ชื่อภาพประกอบ: ผลการทดลองกรณีศึกษาที่ 1]
%     (ก)~การเบี่ยงเบนความถี่ในพื้นที่ 1 (เฮิร์ซ)
%     (ข)~การเบี่ยงเบนความถี่ในพื้นที่ 2 (เฮิร์ซ)}
%   \label{fig:result_case1}
% \end{figure}

%ตาราง~\ref{tab:result_case1} แสดงผลลัพธ์เปรียบเทียบตัวควบคุมในกรณีศึกษาที่ 1

\begin{table}[htbp]
  \centering
  \caption{ชื่อตาราง: ผลลัพธ์การเปรียบเทียบในกรณีศึกษาที่ 1}
  \label{tab:result_case1}
  \begin{tabular}{lllll}
    \toprule
    ชนิดตัวควบคุม & ผลลัพธ์      & $f_1$  & $f_2$  & $f_3$ \\
    \midrule
    \multirow{4}{*}{PID+DD (เสนอ)}
      & Overshoot    & 50.21  & 50.14  & 50.10 \\
      & Undershoot   & 49.88  & 49.95  & 49.97 \\
      & Settling time & 23.34 & 13.82  & 15.69 \\
      & ITAE         & \multicolumn{3}{c}{25.51} \\
    \midrule
    \multirow{4}{*}{PID}
      & Overshoot    & 50.21  & 50.14  & 50.10 \\
      & Undershoot   & 49.86  & 49.93  & 49.96 \\
      & Settling time & 28.56 & 16.30  & 22.37 \\
      & ITAE         & \multicolumn{3}{c}{28.25} \\
    \bottomrule
  \end{tabular}
\end{table}

\subsubsection{กรณีศึกษาที่ 2}

\noindent\hspace{\parindent}%
[อธิบายผลการทดลองในกรณีศึกษาที่ 2]

%%------------------------------------------------------------------
\section{ผลการทดลองส่วนที่ 2}

\subsection{ชื่อหัวข้อย่อย}

\subsubsection{กรณีศึกษาที่ 1}

\noindent\hspace{\parindent}%
[อธิบายผลการทดลองส่วนที่ 2 กรณีศึกษาที่ 1]

\subsubsection{กรณีศึกษาที่ 2}

\noindent\hspace{\parindent}%
[อธิบายผลการทดลองส่วนที่ 2 กรณีศึกษาที่ 2]

%%------------------------------------------------------------------
%% เพิ่มหัวข้อย่อยเพิ่มเติมตามความเหมาะสม
%%------------------------------------------------------------------
